% ================= APPENDIX: CONTROL ALGORITHMS =================
\section{Control Algorithms}
\label{app:control_algorithms}

This appendix holds the implementation detail of the control software whose
architecture is given in Section~\ref{sec:control_sw}. The derivations remain in
Section~\ref{sec:modeling} and are not repeated; what follows is the executable
form --- the loop, the estimator, the control law, the mode logic and the fault
responses --- together with the gains and timing figures that are re-tuned
against hardware and therefore change more often than the design does.

\subsection{Control Loop}
\label{app:loop}

The top-level loop is the fixed sequence of Figure~\ref{fig:software_pipeline},
executed at the sample period $\Delta t$ of
Table~\ref{tab:control2d_vars}. Its ordering is a real-time requirement, not a
convenience: with no dynamic allocation and no blocking call, the worst-case
cycle is bounded by construction.

\begin{verbatim}
every dt:
    t0 <- now()

    # 1. Sensor ingest -- fixed-size, non-blocking reads only
    imu      <- read_imu()                # rates and specific force
    wheels   <- read_driver_replies()     # wheel angle and speed, per axis
    flag_stale(imu, wheels)               # feeds the fault logic below

    # 2. Estimation -- propagate on rates, correct on the absolute reference
    xhat     <- estimator_predict(xhat, imu.rates, dt)
    xhat     <- estimator_correct(xhat, imu.accel)

    # 3. State assembly
    x        <- assemble_state(xhat, wheels)

    # 4. Supervision -- mode first, so the control law below is unambiguous
    mode     <- supervisor_update(mode, x, faults)

    # 5. Control law for the current mode
    u        <- control_law(mode, x, gains[mode])
    u        <- apply_limits(u)           # torque and momentum saturation

    # 6. Output
    send_torque_commands(u)
    set_brake(mode)

    # 7. Telemetry -- side-branch, never in the control path
    telemetry_publish_nonblocking(x, u, mode)

    log_cycle_time(now() - t0)            # feeds the timing budget
\end{verbatim}

Telemetry sits after the output stage and is non-blocking by construction, so
that logging cannot perturb loop timing --- the property that makes the sample
period a guarantee rather than an average.

\subsection{Attitude Estimation}
\label{app:estimator}

The estimator propagates attitude on the measured rates and corrects it against
the absolute reference, while estimating the rate-sensor bias. The two error
sources it must handle pull in opposite directions and are the reason a single
sensor will not do: the rate measurement is clean over a cycle but its bias
drifts, and the absolute reference is drift-free but is corrupted by the
vibration the machine itself generates by spinning the wheels
(Section~\ref{sec:background}).

\begin{verbatim}
estimator_predict(xhat, rates, dt):
    w_corrected <- rates - xhat.bias
    xhat.att    <- integrate_attitude(xhat.att, w_corrected, dt)
    xhat.P      <- propagate_covariance(xhat.P, dt)   # bias random walk
    return xhat

estimator_correct(xhat, accel):
    # Gate the correction: during high-acceleration phases the absolute
    # reference is not measuring the vertical, and applying it would inject
    # the manoeuvre into the attitude estimate.
    if not accel_trustworthy(accel):
        return xhat
    innovation <- reference_direction(accel) - predicted_direction(xhat.att)
    K          <- gain(xhat.P, R_accel)
    xhat       <- apply_correction(xhat, K, innovation)   # attitude and bias
    return xhat
\end{verbatim}

The gating in \texttt{estimator\_correct} is the load-bearing part. The jump-up
and any aggressive recovery are exactly the phases in which the accelerometer
does not measure the vertical, so an ungated correction would drive the estimate
with the manoeuvre. The trust test and its threshold are recorded in
Table~\ref{tab:control_params}.

\subsection{Control Law}
\label{app:control_law}

The balancing law is the state feedback of
Sections~\ref{sec:control2d} and~\ref{sec:3dmodel}, with two additions required
by hardware: the stored-momentum penalty that bleeds off wheel speed, and the
offset filter that absorbs a constant tilt-sensor bias without a calibration
stage.

\begin{verbatim}
control_law(mode, x, K):
    if mode in {EDGE_BALANCE, CORNER_BALANCE, SLEW}:
        e        <- x.att - reference(mode)   # zero for pure regulation
        x_f      <- lowpass(x_f, e)           # absorbs constant sensor offset
        u        <- -K @ [e - x_f, x.rates, x.wheel_speeds]
        return u

    if mode == SPINUP:
        return speed_ramp(x.wheel_speeds, target_speed)

    if mode == JUMPUP:
        return open_loop_jump_sequence(t_since_mode_entry)

    return zero_torque()                      # IDLE and FAULT
\end{verbatim}

The wheel-speed term in the feedback is not merely a damping state. Under any
persistent disturbance the wheels drift toward their speed limit, so the gain on
stored momentum is what keeps the actuators inside their envelope; without it
the machine balances correctly until it saturates and then does not.

\paragraph{Saturation handling.} \texttt{apply\_limits} enforces the torque
limit and the momentum envelope. Torque is clipped per axis; momentum saturation
is handled by the speed penalty above rather than by clipping, since clipping a
saturated wheel removes control authority precisely when it is needed. When a
wheel reaches its limit despite the penalty, the supervisor is notified and
treats it as the fault of Section~\ref{app:faults}.

\subsection{Mode Supervision}
\label{app:modes}

The supervisor implements the transitions of Table~\ref{tab:sw_modes}. Every
threshold carries hysteresis, and every mode has a disarm exit.

\begin{verbatim}
supervisor_update(mode, x, faults):
    if faults.critical:
        return FAULT                      # from any mode, unconditionally

    switch mode:
      IDLE:            if arm_command:            -> SPINUP
      SPINUP:          if wheel_speed >= target:  -> JUMPUP
      JUMPUP:          if |tilt| < capture_band:  -> EDGE_BALANCE
                       if timeout:                -> FAULT
      EDGE_BALANCE:    if |tilt| > release_band:  -> FAULT     # fall detected
                       if corner_command:         -> CORNER_BALANCE
      CORNER_BALANCE:  if |tilt| > release_band:  -> FAULT
                       if slew_command:           -> SLEW
      SLEW:            if reference_reached:      -> CORNER_BALANCE
                       if |tilt| > release_band:  -> FAULT
      FAULT:           if manual_reset and safe:  -> IDLE
    return mode
\end{verbatim}

The capture and release bands differ deliberately --- entry into balancing is
stricter than the condition for staying in it --- which is the hysteresis that
prevents chattering at a mode boundary.

\subsection{Fault Detection and Response}
\label{app:faults}

Table~\ref{tab:fault_responses} gives the detection test and response for each
failure case. The distinction that matters is between faults the machine can
continue through and faults it cannot: only the loss of the balancing
equilibrium terminates the run, because a body outside its capture region cannot
be recovered by the linear law, and continued actuation adds energy to a fall
that is already happening.

\begin{table}[h]
\centering
\caption{Fault detection and response. Only loss of equilibrium terminates the
run; the remaining cases keep the machine controlled.}
\label{tab:fault_responses}
\small
\begin{tabular}{L{3.2cm} L{4.6cm} L{5.4cm}}
\toprule
Fault & Detection & Response \\
\midrule
Sensor dropout / stale data & Missing or repeated sample within the cycle deadline & Propagate on the remaining sensor; escalate to FAULT if persistent beyond the tolerance \\
Bus error       & Error-frame count or missed driver reply & Retry within the cycle; escalate on repetition \\
Encoder error flag & Flag set, or angle discontinuity beyond the physical rate limit & Disarm the affected axis; FAULT if it is load-bearing for the current mode \\
Wheel saturation & Speed at limit despite the momentum penalty & Notify supervisor; reduce reference authority; FAULT if sustained \\
Loss of equilibrium & Tilt outside the release band & FAULT: wheels commanded down, drivers disarmed \\
Loop overrun    & Measured cycle time exceeds the deadline & Log; FAULT if repeated, since the timing guarantee no longer holds \\
Watchdog        & Loop fails to service the timer & Hardware disarm, independent of the software path \\
\bottomrule
\end{tabular}
\end{table}

\subsection{Timing Budget}
\label{app:timing}

Table~\ref{tab:timing_budget} records the per-stage timing budget and the
measured margin. It is re-measured whenever a stage changes, because the
guarantee of Section~\ref{sec:control_sw} is a claim about the worst case and
not about the mean.

\begin{table}[h]
\centering
\caption{Per-cycle timing budget against the sample period. Re-measured
whenever a stage changes.}
\label{tab:timing_budget}
\small
\begin{tabular}{L{4.6cm} c c}
\toprule
Stage & Budget & Measured worst case \\
\midrule
Sensor ingest        & \textcolor{red}{TBD} & \textcolor{red}{TBD} \\
Estimation           & \textcolor{red}{TBD} & \textcolor{red}{TBD} \\
State assembly       & \textcolor{red}{TBD} & \textcolor{red}{TBD} \\
Supervision          & \textcolor{red}{TBD} & \textcolor{red}{TBD} \\
Control law          & \textcolor{red}{TBD} & \textcolor{red}{TBD} \\
Command output       & \textcolor{red}{TBD} & \textcolor{red}{TBD} \\
\midrule
\textbf{Total cycle} & \textcolor{red}{\textbf{TBD}} & \textcolor{red}{\textbf{TBD}} \\
\textbf{Margin to deadline} & --- & \textcolor{red}{\textbf{TBD}} \\
\bottomrule
\end{tabular}
\end{table}

\subsection{Gains and Tuning Parameters}
\label{app:gains}

Table~\ref{tab:control_params} is the tuning record: every constant the control
software carries, with the source that fixes it. Values computed from the model
are distinguished from values tuned against hardware, since only the latter must
be re-established when the machine is re-built.

\begin{table}[h]
\centering
\caption{Control and estimation parameters. ``Model'' values follow from the
design; ``bench'' values are tuned against hardware and re-established after a
re-build.}
\label{tab:control_params}
\small
\begin{tabular}{l L{5.4cm} c c}
\toprule
Symbol & Parameter & Source & Value \\
\midrule
$\Delta t$        & Sample period                      & Model & \textcolor{red}{TBD} \\
$Q, R$            & State and input weights            & Model & \textcolor{red}{TBD} \\
$K_d$ (edge)      & Feedback gain, edge balance        & Model & \textcolor{red}{TBD} \\
$K_d$ (corner)    & Feedback gain, corner balance      & Model & \textcolor{red}{TBD} \\
---               & Momentum-penalty weight            & Bench & \textcolor{red}{TBD} \\
$\gamma$          & Estimator correction weight        & Bench & \textcolor{red}{TBD} \\
---               & Accelerometer trust threshold      & Bench & \textcolor{red}{TBD} \\
$\alpha$          & Offset-filter coefficient          & Bench & \textcolor{red}{TBD} \\
$\theta_{\text{th}}$ & Balancing capture band          & Bench & \textcolor{red}{TBD} \\
---               & Balancing release band (hysteresis)& Bench & \textcolor{red}{TBD} \\
---               & Spin-up target wheel speed         & Model & \textcolor{red}{TBD} \\
---               & Brake actuation timing             & Bench & \textcolor{red}{TBD} \\
---               & Torque limit per axis              & Model & \textcolor{red}{TBD} \\
---               & Watchdog timeout                   & Model & \textcolor{red}{TBD} \\
\bottomrule
\end{tabular}
\end{table}
